\chapter{Continue toevalsvariabelen}\label{ch:g12:contdist}

Wachttijden, natuurkundige metingen, aandelen: veel toevalsgrootheden nemen
een continuüm aan waarden aan, en geen enkele afzonderlijke waarde heeft een
\omterm{def:g10:proba:distribution}{kans} verschillend van nul. Kansen worden dan berekend door een
\emph{\omterm{def:g12:contdist:density}{kansdichtheid}} te integreren. Dit hoofdstuk bestudeert de uniforme, de
exponentiële en de \omterm{def:g12:contdist:normal}{normale verdeling}, en gebruikt de \omterm{def:g12:contdist:normal}{normale verdeling} om de
schommelingen van peilingen en steekproeven te meten.

\section{Kansdichtheden}

\begin{definition}[Kansdichtheid, continue toevalsvariabele]\label{def:g12:contdist:density}
Een \emph{kansdichtheid}\index{kansdichtheid} op een \omterm{def:g10:numbers:interval}{interval} $I$ is een
continue, niet-negatieve \omterm{def:g11:func:function}{functie} $f$ op $I$ met $\int_I f(t)\,\dd t = 1$
(waarbij de \omterm{def:g12:integ:area}{integraal} over een onbegrensde $I$ opgevat wordt als een limiet
van integralen over groeiende begrensde intervallen). Een \omterm{def:g12:randvar:rv}{toevalsvariabele} $X$
heeft dichtheid $f$ als voor alle $a \leq b$ in $I$ geldt
\[
\P(a \leq X \leq b) = \int_a^b f(t)\,\dd t .
\]
\end{definition}

Kansen zijn \emph{oppervlaktes onder de kromme van de dichtheid}. In het
bijzonder is $\P(X = a) = \int_a^a f = 0$ voor elke afzonderlijke waarde $a$:
alleen intervallen dragen \omterm{def:g10:proba:distribution}{kans}, en $\P(a \leq X \leq b) = \P(a < X < b)$.

\begin{definition}[Verwachtingswaarde en variantie]\label{def:g12:contdist:exp}
Voor $X$ met dichtheid $f$ op $I$:
\[
\E(X) = \int_I t\,f(t)\,\dd t, \qquad
\V(X) = \int_I \bigl(t - \E(X)\bigr)^2 f(t)\,\dd t
= \E(X^2) - \E(X)^2 .
\]
\end{definition}

De regels van de hoofdstukken 14 en 15 (de lineariteit van de
\omterm{def:g12:randvar:exp}{verwachtingswaarde}, de optelbaarheid van de \omterm{def:g12:randvar:exp}{variantie} voor onafhankelijke
variabelen, Bienaymé--Chebyshev) blijven geldig; hun bewijzen, met integralen
in plaats van sommen, worden op dit niveau aangenomen.

\section{Uniforme verdeling}

\begin{definition}[Uniforme verdeling]\label{def:g12:contdist:uniform}
$X$ volgt de \emph{uniforme verdeling}\index{uniforme verdeling}
$\mathcal U\intcc ab$ als haar dichtheid constant is,
$f(t) = \frac{1}{b-a}$ op $\intcc{a}{b}$. Dan geldt voor
$\intcc{c}{d} \subseteq \intcc ab$ dat
$\P(c \leq X \leq d) = \frac{d - c}{b - a}$: de \omterm{def:g10:proba:distribution}{kans} is evenredig met de
lengte.
\end{definition}

\begin{proposition}\label{prop:g12:contdist:uniformmoments}
Is $X \sim \mathcal U\intcc ab$, dan is
$\E(X) = \dfrac{a+b}{2}$ en $\V(X) = \dfrac{(b-a)^2}{12}$.
\end{proposition}

\begin{proof}
$\E(X) = \int_a^b \frac{t}{b-a}\dd t = \frac{b^2 - a^2}{2(b-a)} =
\frac{a+b}{2}$. Voor de \omterm{def:g12:randvar:exp}{variantie} is
$\E(X^2) = \frac{b^3 - a^3}{3(b-a)} = \frac{a^2 + ab + b^2}{3}$, en
\[
\V(X) = \frac{a^2 + ab + b^2}{3} - \frac{(a+b)^2}{4}
= \frac{4a^2 + 4ab + 4b^2 - 3a^2 - 6ab - 3b^2}{12}
= \frac{(b - a)^2}{12}. \qedhere
\]
\end{proof}

\section{Exponentiële verdeling}

\begin{definition}[Exponentiële verdeling]\label{def:g12:contdist:expo}
Voor $\lambda > 0$ volgt $X$ de \emph{exponentiële
verdeling}\index{exponentiële verdeling} $\mathcal E(\lambda)$ als haar
dichtheid op $\intco{0}{+\infty}$ gegeven wordt door
\[
f(t) = \lambda\,\eu^{-\lambda t}.
\]
Dan is $\P(X \leq x) = 1 - \eu^{-\lambda x}$ en
$\P(X > x) = \eu^{-\lambda x}$ voor $x \geq 0$.
\end{definition}

\begin{proposition}\label{prop:g12:contdist:expomoments}
Is $X \sim \mathcal E(\lambda)$, dan is $\E(X) = \dfrac{1}{\lambda}$ en
$\V(X) = \dfrac{1}{\lambda^2}$.
\end{proposition}

\begin{proof}
\omterm{thm:g12:integ:ibp}{Partiële integratie} op $\intcc{0}{A}$:
\[
\int_0^A t\,\lambda\eu^{-\lambda t}\dd t
= \bigl[-t\,\eu^{-\lambda t}\bigr]_0^A + \int_0^A \eu^{-\lambda t}\dd t
= -A\eu^{-\lambda A} + \frac{1 - \eu^{-\lambda A}}{\lambda}
\xrightarrow[A\to+\infty]{} \frac1\lambda,
\]
waarbij $A\eu^{-\lambda A} \to 0$ (\cref{thm:g12:exp:growth}). Een tweede
\omterm{thm:g12:integ:ibp}{partiële integratie} geeft $\E(X^2) = \frac{2}{\lambda^2}$, waaruit
$\V(X) = \frac{2}{\lambda^2} - \frac{1}{\lambda^2}$.
\end{proof}

\begin{theorem}[Geheugenloosheid]\label{thm:g12:contdist:memoryless}
\index{geheugenloosheid}
Is $X \sim \mathcal E(\lambda)$, dan geldt voor alle $s, t \geq 0$:
\[
\pcond{X > s}{X > s + t} = \P(X > t) .
\]
De \omterm{def:g12:contdist:expo}{exponentiële verdeling} is de verdeling van levensduren \emph{zonder
veroudering} (radioactieve kernen, geen gloeilampen).
\end{theorem}

\begin{proof}
\[
\pcond{X > s}{X > s+t}
= \frac{\P(X > s + t)}{\P(X > s)}
= \frac{\eu^{-\lambda(s+t)}}{\eu^{-\lambda s}} = \eu^{-\lambda t}
= \P(X > t). \qedhere
\]
\end{proof}

\begin{omfigure}
\begin{tikzpicture}
\begin{axis}[omaxis,
  width=8.8cm, height=5cm,
  xmin=-0.4, xmax=4.4, ymin=0, ymax=1.15,
  xtick={1.3}, xticklabels={$t$}, ytick={1}, yticklabels={$\lambda$}]
  \addplot[name path=dens, omDef, thick, domain=0:4.2, samples=80]
    {exp(-x)};
  \path[name path=xaxis] (0,0) -- (4.2,0);
  \addplot[omThm!20] fill between[of=dens and xaxis,
    soft clip={domain=1.3:4.2}];
  \draw[black, dashed, thin] (axis cs:1.3,0) -- (axis cs:1.3,0.2725);
  \node[omThm, font=\small, anchor=west] at (axis cs:2.1,0.35)
    {$\P(X > t) = \eu^{-\lambda t}$};
\end{axis}
\end{tikzpicture}

{\small De exponentiële dichtheid $\lambda\eu^{-\lambda x}$ (hier
$\lambda = 1$): de oppervlakte van de staart voorbij $t$ (rood) is
$\eu^{-\lambda t}$, en de geheugenloosheid zegt dat elke staart eruitziet als
de hele verdeling, herschaald.}
\end{omfigure}

\section{Normale verdeling}

\begin{definition}[Normale verdeling]\label{def:g12:contdist:normal}
$X$ volgt de \emph{standaardnormale verdeling}\index{normale verdeling}
$\mathcal N(0, 1)$ als haar dichtheid op $\R$ gegeven wordt door
\[
\varphi(t) = \frac{1}{\sqrt{2\pi}}\,\eu^{-t^2/2}
\]
(de klokkromme). Algemener is $X \sim \mathcal N(\mu, \sigma^2)$ als
$\dfrac{X - \mu}{\sigma} \sim \mathcal N(0,1)$; dan is $\E(X) = \mu$ en
$\V(X) = \sigma^2$.
\end{definition}

\begin{omfigure}
\begin{tikzpicture}
\begin{axis}[
  width=11cm, height=4.6cm,
  axis lines=middle, xmin=-4, xmax=4, ymin=0, ymax=0.45,
  xtick={-3,-2,-1,0,1,2,3},
  xticklabels={$\mu{-}3\sigma$,$\mu{-}2\sigma$,$\mu{-}\sigma$,$\mu$,
               $\mu{+}\sigma$,$\mu{+}2\sigma$,$\mu{+}3\sigma$},
  ytick=\empty, samples=120, domain=-4:4]
  \addplot[name path=gauss, omDef, thick] {exp(-x^2/2)/sqrt(2*pi)};
  \path[name path=xaxis] (-4,0) -- (4,0);
  \addplot[omDef!20] fill between[of=gauss and xaxis,
           soft clip={domain=-2:2}];
  \node[omDef] at (axis cs:0,0.17) {$\approx 95.4\%$};
\end{axis}
\end{tikzpicture}

{\small De normale dichtheid: ongeveer $68.3\%$ van de massa ligt binnen
$\sigma$ van het \omterm{def:g11:stat:mean}{gemiddelde}, $95.4\%$ binnen $2\sigma$ en $99.7\%$ binnen
$3\sigma$.}
\end{omfigure}

\begin{remark}
De factor $\frac{1}{\sqrt{2\pi}}$ maakt de totale oppervlakte $1$ --- een
beroemde berekening (de \emph{\omterm{def:g12:integ:area}{integraal} van Gauss}) die aan de universiteit
gemaakt wordt. Er is geen elementaire formule voor $\int \varphi$: normale
kansen lees je af uit tabellen of van een rekentoestel.
\end{remark}

\begin{theorem}[De Moivre--Laplace, aangenomen]\label{thm:g12:contdist:tcl}
\index{stelling van De Moivre--Laplace}
Zij $X_n \sim \mathcal B(n, p)$ en $Z_n = \dfrac{X_n - np}{\sqrt{np(1-p)}}$
(de gestandaardiseerde binomiale variabele). Dan geldt voor alle
$a \leq b$:
\[
\P(a \leq Z_n \leq b) \xrightarrow[n\to+\infty]{}
\int_a^b \varphi(t)\,\dd t .
\]
\emph{Dit resultaat wordt op dit niveau aangenomen.} Het verklaart waarom de
klokkromme overal opduikt: een \omterm{def:g12:randvar:binomial}{binomiale verdeling} met grote $n$ is bij
\omterm{def:g10:numbers:approx}{benadering} normaal --- en (centrale limietstelling, universiteit) elke som van
veel kleine onafhankelijke effecten eveneens.
\end{theorem}

\begin{method}[Schommelingsinterval op $95\%$]\label{met:g12:contdist:fluctuation}
\index{betrouwbaarheidsinterval}
Omdat een normale variabele met \omterm{def:g10:proba:distribution}{kans} $0.95$ binnen $1.96$
\omterm{def:g11:stat:variance}{standaardafwijkingen} van haar \omterm{def:g11:stat:mean}{gemiddelde} valt, voldoet voor grote $n$ een
\omterm{def:g10:stats:series}{frequentie} $F_n = \frac{X_n}{n}$ uit een \omterm{def:g10:proba:sample}{steekproef} $\mathcal B(n,p)$ aan
\[
\P\left(p - 1.96\sqrt{\tfrac{p(1-p)}{n}} \leq F_n \leq
p + 1.96\sqrt{\tfrac{p(1-p)}{n}}\right) \approx 0.95 .
\]
\begin{itemize}
  \item \emph{Schommelingsinterval} (toetsen): legt een beweerde $p$ de
        waargenomen \omterm{def:g10:stats:series}{frequentie} buiten dit \omterm{def:g10:numbers:interval}{interval}, verwerp dan de bewering
        op het niveau $5\%$.
  \item \emph{\omterm{met:g12:contdist:fluctuation}{Betrouwbaarheidsinterval}} (schatten): het vereenvoudigde
        \omterm{def:g10:numbers:interval}{interval}
        $\left[F_n - \frac{1}{\sqrt n},\ F_n + \frac{1}{\sqrt n}\right]$
        bevat $p$ met \omterm{def:g10:proba:distribution}{kans} minstens $0.95$ (met
        $1.96\sqrt{p(1-p)} \leq 1$, zie \cref{exo:g12:sums:8}).
\end{itemize}
\end{method}

\begin{example}\label{ex:g12:contdist:poll}
Een peiling bij $n = 1000$ kiezers geeft een kandidaat $52\%$. Het
\omterm{met:g12:contdist:fluctuation}{betrouwbaarheidsinterval} $52\% \pm \frac{1}{\sqrt{1000}} \approx 52\% \pm
3.2\%$ bevat nog altijd $50\%$: de peiling alleen toont niet aan dat de
kandidaat voorstaat.
\end{example}

\section{Oefeningen}

\begin{exercise}[$\star$]\label{exo:g12:contdist:1}
Er passeert om de $15$ minuten een bus; een reiziger komt op een uniform
willekeurig tijdstip aan. Zij $X \sim \mathcal U\intcc{0}{15}$ de wachttijd.
Bereken $\P(X \leq 5)$, $\P(4 \leq X \leq 10)$, $\E(X)$ en $\sigma(X)$.
\end{exercise}

\begin{exercise}[$\star$]\label{exo:g12:contdist:2}
Ga na dat $f(t) = \frac{3}{8}t^2$ een dichtheid is op $\intcc{0}{2}$, en
bereken $\E(X)$ en $\P(X \geq 1)$ voor een variabele $X$ met die dichtheid.
\end{exercise}

\begin{exercise}[$\star$]\label{exo:g12:contdist:3}
De levensduur (in jaren) van een elektronisch onderdeel volgt
$\mathcal E(0.2)$.
\begin{enumerate}
  \item Bereken de verwachte levensduur en $\P(X > 5)$.
  \item Het onderdeel heeft al $3$ jaar gewerkt. Wat is de \omterm{def:g10:proba:distribution}{kans} dat het nog
        minstens $5$ jaar werkt?
\end{enumerate}
\end{exercise}

\begin{exercise}[$\star\star$]\label{exo:g12:contdist:4}
De halveringstijd van een radioactief atoom met een levensduur die
$\mathcal E(\lambda)$ volgt, is de \omterm{def:g11:stat:median}{mediaan} $m$: $\P(X > m) = \frac12$. Druk
$m$ uit in \omterm{def:g11:func:function}{functie} van $\lambda$ en vergelijk met de \omterm{def:g12:randvar:exp}{verwachtingswaarde}. Welke
van beide is groter, en waarom is dat zinvol voor een scheve verdeling?
\end{exercise}

\begin{exercise}[$\star\star$]\label{exo:g12:contdist:5}
De lichaamslengtes in een bevolking volgen $\mathcal N(175,\ 7^2)$ (in cm).
Schat met de regel $68$--$95$--$99.7$ het aandeel van de bevolking met een
lengte tussen $168$ en $182$ cm, boven $189$ cm, en onder $154$ cm.
\end{exercise}

\begin{exercise}[$\star\star$]\label{exo:g12:contdist:6}
Een machine vult zakken met het opschrift $500$ g; de gevulde massa volgt
$\mathcal N(\mu,\ 4^2)$. De regelgeving eist dat hoogstens $2.5\%$ van de
zakken minder dan $500$ g weegt. Welke minimale instelling van $\mu$ voldoet,
volgens de regel van $95\%$?
\end{exercise}

\begin{exercise}[$\star\star$]\label{exo:g12:contdist:7}
Een dobbelsteen die zuiver heet te zijn, wordt $1200$ keer geworpen en toont
$260$ keer een zes (zoals in \cref{exo:g12:sums:7}).
\begin{enumerate}
  \item Bereken het schommelingsinterval op $95\%$ voor de \omterm{def:g10:stats:series}{frequentie} van
        zessen van een zuivere dobbelsteen over $1200$ worpen.
  \item Ligt de waargenomen \omterm{def:g10:stats:series}{frequentie} erbinnen? Vergelijk de kracht van dit
        besluit met de analyse via Bienaymé--Chebyshev.
\end{enumerate}
\end{exercise}

\begin{exercise}[$\star\star\star$]\label{exo:g12:contdist:8}
Vóór een verkiezing zal een peiling bij $n$ mensen de score $p$ van een
kandidaat schatten met de waargenomen \omterm{def:g10:stats:series}{frequentie} $F_n$.
\begin{enumerate}
  \item Welke omvang van de \omterm{def:g10:proba:sample}{steekproef} waarborgt met het
        \omterm{met:g12:contdist:fluctuation}{betrouwbaarheidsinterval} van \cref{met:g12:contdist:fluctuation} een
        marge van $\pm 2$ punten?
  \item In werkelijkheid liggen de kandidaten $1$ punt uit elkaar. Leg uit
        waarom geen enkele realistische peiling de winnaar betrouwbaar kan
        aanwijzen, hoe goed ze ook uitgevoerd wordt.
\end{enumerate}
\end{exercise}

\begin{exercise}[$\star\star\star$]\label{exo:g12:contdist:9}
Zij $X$ een variabele met dichtheid $f$ op $\intcc ab$ en zij
$Y = \alpha X + \beta$ met $\alpha > 0$.
\begin{enumerate}
  \item Toon aan dat $\P(c \leq Y \leq d)
        = \P\left(\frac{c - \beta}{\alpha} \leq X \leq
        \frac{d-\beta}{\alpha}\right)$, en leid af dat $Y$ als dichtheid
        $g(y) = \frac{1}{\alpha} f\left(\frac{y - \beta}{\alpha}\right)$
        heeft.
  \item Leid af dat $\mu + \sigma Z$ voor $Z \sim \mathcal N(0,1)$ als
        dichtheid $\frac{1}{\sigma\sqrt{2\pi}}\,
        \eu^{-(y - \mu)^2/(2\sigma^2)}$ heeft --- de algemene normale
        dichtheid.
\end{enumerate}
\end{exercise}

\section{Opgave: wachten op de bus, de menigte meten}

\begin{problem}[{Weekendopgave --- drie dichtheden besturen de wereld: de
volslagen onwetendheid, het geheugenloze wachten en de klok van vele kleine
oorzaken; plus de paradox die je bus altijd te laat maakt}]
\label{pb:g12:contdist:1}
Waarom lijkt je bus altijd langer op zich te laten wachten dan het uurrooster
belooft? Waarom zit een \omterm{def:g11:stat:mean}{gemiddelde} leerling in een grotere klas dan de
\omterm{def:g11:stat:mean}{gemiddelde} klas? Waarom hebben je vrienden meer vrienden dan jij? Eén sluw
stukje wiskunde --- het naar lengte vertekende bemonsteren --- beantwoordt alle
drie de vragen, en het woont in dit hoofdstuk, tussen de uniforme, de
exponentiële en de \omterm{def:g12:contdist:normal}{normale verdeling} (\cref{def:g12:contdist:uniform},
\cref{def:g12:contdist:expo}, \cref{def:g12:contdist:normal}). Deze laatste
opgave van het volume zet de drie dichtheden aan het werk en eindigt waar de
hele reeks naar wijst: bij de klokkromme.

\medskip
\noindent\textbf{Deel I --- Drie personages.}
\begin{enumerate}
  \item Een bus komt op een uniform willekeurig tijdstip in de komende $30$
        minuten aan. Bereken $\P(\text{wachttijd} \leq 10)$, de verwachte
        wachttijd en $\sigma$ (\cref{prop:g12:contdist:uniformmoments}).
  \item Ga na dat $f(x) = 2x$ op $\intcc{0}{1}$ een dichtheid is, en bereken
        $\E(X)$ en $\P\left(X > \frac12\right)$.
  \item Telefoonoproepen bereiken een hulplijn als een geheugenloze stroom: de
        wachttijd tot de volgende oproep is exponentieel met \omterm{def:g11:stat:mean}{gemiddelde} $10$
        minuten. Bereken $\P(T > 15)$ en de mediane wachttijd --- waarom is de
        \omterm{def:g11:stat:median}{mediaan} \emph{kleiner} dan het \omterm{def:g11:stat:mean}{gemiddelde}?
  \item Je hebt al $20$ minuten gewacht. Wat is $\P(T > 35 \mid T > 20)$
        (\cref{thm:g12:contdist:memoryless})? Duid het in één zin.
  \item Koppel elke wachttijd aan haar juiste model --- uniform, exponentieel
        of geen van beide: (a) de klik van een geigerteller; (b) een metro die
        om de $8$ minuten rijdt, waarbij je aankomst niet gelijkloopt; (c) de
        volgende oproep op de hulplijn; (d) het doorbranden van een gloeilamp
        die \emph{verslijt}. Verantwoord (d) met de stelling over de
        geheugenloosheid.
\end{enumerate}

\medskip
\noindent\textbf{Deel II --- De klok aan het werk.}
(Gebruik de normale ijkpunten: $68\,\%$ binnen $\sigma$, $95\,\%$ binnen
$2\sigma$, $99.7\,\%$ binnen $3\sigma$.)
\begin{enumerate}[resume]
  \item De lengtes van volwassenen volgen ruwweg $\mathcal N(172,\ 8)$ (cm).
        Welke aandelen liggen in $\intcc{164}{180}$, $\intcc{156}{188}$ en
        $\intcc{148}{196}$?
  \item Schat het aandeel groter dan $190$ cm (eerst de \omterm{pb:g11:stat:1}{z-score}; daarna tabel
        of rekentoestel).
  \item IQ-scores worden geijkt op $\mathcal N(100, 15)$. Welk aandeel scoort
        boven $130$ --- en ruwweg één persoon op hoeveel?
  \item Een machine snijdt bouten met lengtes $\mathcal N(50,\ 0.1)$ mm en de
        specificatie is $\intcc{49.8}{50.2}$. Welk aandeel wordt afgekeurd? De
        industrie viert processen ``op zes sigma'', waarvan de specificaties
        op $\pm 6\sigma$ liggen: wat levert dat op, en waarom jagen fabrieken
        erop?
  \item De Moivre--Laplace (\cref{thm:g12:contdist:tcl}): benader voor $100$
        worpen met een zuivere munt $\P(45 \leq \text{aantal kruis} \leq 55)$
        (continuïteitscorrectie: $\pm 0.5$; $\sigma = 5$). Welke houten
        machine uit \cref{pb:g11:binom:1} maakt deze stelling glad tot een
        kromme?
  \item Uit \cref{exo:g12:contdist:8}: een marge van $\pm 2$ punten vraagt
        ongeveer $n = 2\,400$. Bereken de \omterm{def:g10:proba:sample}{steekproef} die nodig is voor
        $\pm 0.5$ punten --- en besluit, in één zin, waarom geen enkele
        peiling een race met één punt \omterm{def:g11:seq:arithmetic}{verschil} eerlijk kan beslechten.
\end{enumerate}

\medskip
\noindent\textbf{Deel III --- De busparadox.}
\begin{enumerate}[resume]
  \item Bussen komen aan als een geheugenloze stroom met \omterm{def:g11:stat:mean}{gemiddelde} tussentijd
        $10$ minuten. Je bereikt de halte op een willekeurig ogenblik. De
        intuïtie zegt dat de verwachte wachttijd $5$ minuten is (een halve
        tussentijd). Wat zegt de stelling over de geheugenloosheid in de
        plaats?
  \item Stel de intuïtie op de proef met een speelgoeduurrooster: de
        tussentijden wisselen af tussen $5$ en $15$ minuten (\omterm{def:g11:stat:mean}{gemiddelde}
        tussentijd $10$). Bereken de \omterm{def:g10:proba:distribution}{kans} dat een uniform willekeurige aankomst
        in een lange tussentijd valt, en daarna de echte verwachte wachttijd
        --- en noem de schuldige: aankomsten bemonsteren de tussentijden met
        een \omterm{def:g10:proba:distribution}{kans} evenredig met hun \emph{lengte}.
  \item In de geheugenloze stroom heeft de tussentijd waarin je terechtkomt
        een verwachte lengte van $20$ minuten --- het dubbele van de typische
        tussentijd (de tijd terug naar de vorige bus en die vooruit naar de
        volgende zijn \emph{allebei} exponentieel met \omterm{def:g11:stat:mean}{gemiddelde} $10$). Breng
        dit in overeenstemming met vraag 12 en formuleer de
        \emph{inspectieparadox} in één zin.
  \item Dezelfde angel aan land: een universiteit heeft negen klassen van $10$
        studenten en één van $110$. Bereken de \omterm{def:g11:stat:mean}{gemiddelde} klasgrootte, en
        daarna de klasgrootte die de \emph{\omterm{def:g11:stat:mean}{gemiddelde} student} beleeft. Welk
        getal zal de brochure vermelden, en welk getal is de geleefde
        werkelijkheid van de studenten?
  \item De vriendschapsparadox --- ``je vrienden hebben gemiddeld meer
        vrienden dan jij'' --- is dezelfde wiskunde. Zeg in één zin wat er de
        rol van de lange tussentijd speelt.
\end{enumerate}

\medskip
\noindent\textbf{Deel IV --- Het volume afsluiten.}
\begin{enumerate}[resume]
  \item Simulatie, de brug van de praktijk: is $U$ uniform op
        $\intoo{0}{1}$, toon dan aan dat $T = -10\ln U$ exponentieel is met
        \omterm{def:g11:stat:mean}{gemiddelde} $10$ (bereken $\P(T > t)$). Elke simulatie van willekeurig
        wachten in de industrie draait op deze truc van één regel.
  \item Een oud recept van spelprogrammeurs bouwt een klok: tel twaalf
        onafhankelijke uniforme variabelen op $\intcc{0}{1}$ op en trek er $6$
        van af. Geef het \omterm{def:g11:stat:mean}{gemiddelde} en de \omterm{def:g12:randvar:exp}{variantie} van de som, en leg uit ---
        met de bord van Galton in de hand --- waarom het resultaat bijna
        normaal is.
  \item Markten storten harder in dan de klok toelaat: de crash van 1987 werd
        omschreven als een ``\omterm{def:g11:prob:model}{gebeurtenis} op $25\sigma$'', wat onder normaliteit
        een \omterm{def:g10:proba:distribution}{kans} van ongeveer $10^{-137}$ heeft. Wat is het juiste besluit ---
        over de wereld, of over het model? (Eén zin, de catechismus van de
        statisticus.)
  \item Slotstuk, en afscheid van het volume: de drie personages in telkens
        één regel --- uniform (onwetendheid binnen grenzen), exponentieel
        (risico zonder geheugen), normaal (de som van vele kleine oorzaken);
        de inspectieparadox als de angel van het hoofdstuk; en de boog van de
        hele reeks --- van het tellen van pakjes sap tot de klokkromme die
        elke menigte meet --- met de universitaire volumes die wachten waar de
        limieten, de integralen en de centrale limietstelling het overnemen.
\end{enumerate}
\end{problem}
